% =====================================================================
%  Harald Høffding, LEDENDE TANKER I DET NITTENDE AARHUNDREDE.
%  København: Gyldendalske Boghandel, Nordisk Forlag, 1920.  Antiqua.
%  (Swedish edition: »Ledande tankar under det nittonde århundradet«, Stockholm: Norstedt, 1921, 146 pp.)
%
%  STATUS: IN PROGRESS (2026-10-23).  An EXTRACT: printed pp. 28-32 only, from the end of the section on Leibniz and
%  Robert Mayer to the opening of »b. Udviklingslæren«.  The text begins and ends mid-sentence.  No title page, no
%  front matter, no Indhold yet; the heading of section »a.« (before p. 28) is not in hand.  Not to be marked complete.
%
%  EDITION.  The text of record is Hans Halvorson's own copy, photographed with Genius Scan (2026-10-23), 5 pages,
%  150-190 ppi greyscale JPEG with Genius Scan's OCR layer: ~/bibliotek/Høffding, Harald/1920-ledende-tanker-pp28-32.pdf
%  (SCANS.tsv by sha256).  PDF 1-5 = pp. 28-32.  The shots curve upward toward the outer margin, so the right-hand words
%  of each line are blurred and slanted; the inner part of every line is sharp.  No e-book, no other scan and no earlier
%  transcription found (checked 2026-10-22/23: the repo, ~/bibliotek, Internet Archive, HathiTrust, the web; Google
%  Books and KB could not be searched by machine).  A modern print-on-demand facsimile exists (ISBN 9781013177927).
%
%  ACCURACY (2026-10-23).  Following the Rydberg-review precedent for a piece this short (TRANSCRIPTION-PLAYBOOK.md §6):
%  every line read at the image in the main session, in horizontal bands at native resolution (15 bands, the 2-3x zoom
%  used for the two misprints below), with two machine readings at hand: tesseract 5 (dan, tessdata_best, psm 6) and
%  Genius Scan's own layer.  The reading was then voted word by word against both machines (punctuation and quotation
%  marks ignored by the vote; those were settled at the image only): in ~1,950 words, no place where both machines
%  agreed against the reading other than digit/case noise at the curled edge (»1« for »i«, »0« for »o.«, »Søger«), and
%  no word the machines had that the reading lacked.  The machines disagree with each other on 72% of lines (149 of 206), almost
%  all at the right edge, so they were pointers only, never copy-text.  OCRDIFF: tesseract+layer | 104 candidates |
%  0 corrected | all witness's fault | 0 unresolved.
%
%  CONVENTIONS.  Diplomatic: orthography, »Mayer’s«/»Carnot’s«, „…“ quotation marks and the dash »—« as printed.
%  Words divided at a line end are joined; a word divided across a page turn carries the marker inside it
%  (Sæt\opage{29}ningen, opbyg\opage{32}ges).  Running heads are recorded in comments at each page turn.  The section
%  heading »b. Udviklingslæren.« is set as printed: »b.« roman, the title italic, centred.
%
%  PRINTER'S ERRORS (transcribed as printed, logged at the site):
%    p. 31  »Zeemann« for Zeeman (in the quotation from C. Christiansen).
%    p. 31  »Mécanique Céliste« for Mécanique céleste (in the quotation from Woodward).
%  Quote balance: p. 31 opens „Principia“ inside an open „…“ quotation (Woodward), as printed.
% =====================================================================
\documentclass[12pt,a4paper,openany]{book}
\usepackage[utf8]{inputenc}
\usepackage[T1]{fontenc}
\usepackage{libertinus}
\usepackage{libertinust1math}
\usepackage[danish]{babel}
\usepackage[protrusion=true,expansion=true]{microtype}
\usepackage{setspace}
\usepackage[margin=1.3in]{geometry}
\usepackage{fancyhdr}
\usepackage{hyperref}
\hypersetup{pdftitle={Ledende Tanker i det nittende Aarhundrede (pp. 28--32)}, pdfauthor={Harald Høffding}, hidelinks}
\pagestyle{fancy}
\fancyhf{}
\fancyhead[LE,RO]{\thepage}
\fancyhead[C]{\textit{Høffding: Ledende Tanker i det nittende Aarhundrede (1920)}}
\setstretch{1.3}
\usepackage{marginnote}
\newcommand{\opage}[1]{\marginnote{\footnotesize\textit{[#1]}}}
\newcommand{\sectionhead}[2]{\par\bigskip\phantomsection\begin{center}#1 \textit{#2}\end{center}\medskip\ignorespaces}

\begin{document}

\begin{center}
{\Large Harald Høffding}\\[0.6em]
{\LARGE Ledende Tanker i det nittende Aarhundrede}\\[0.6em]
København: Gyldendalske Boghandel, Nordisk Forlag, 1920\\[1.2em]
{\small Extract: pp.~28--32. The text begins and ends mid-sentence.}
\end{center}
\bigskip

\noindent[\dots]
% --- p. 28 ---
% running head: LEIBNIZ OG ROBERT MAYER
\opage{28}til Et med Forholdet mellem Grund og Følge, skal her ikke drøftes.
Men for en historisk Betragtning er det af stor Interesse at se, at
Opdagelsen af Energiens Bestaaen var en Videreførelse og Erfaringsbekræftelse
af en Sætning, der allerede var opstillet i det 17. Aarhundrede.
Descartes opstillede en Lov om Bevægelsens Bestaaen,
ifølge hvilken ingen Bevægelse blev til intet, og ingen Bevægelse opstod
af intet; der var i Universet stedse den samme Sum af Bevægelse,
kun fordeltes den paa forskellig Maade. Leibniz kritiserede
denne Lære, idet han hævdede, at ifølge den vilde enhver Overgang
fra Bevægelse til Hvile og fra Hvile til Bevægelse indenfor Universet
være ubegribelig. Der er, hævder han fremdeles, ingen absolut
Modsætning mellem Hvile og Bevægelse; Hvile er stanset Bevægelse,
der skyldes Ligevægt mellem Kræfter (Tendenser), og Bevægelsen
udløses paany, saasnart Ligevægten ophæves. Det er da ikke Bevægelsen,
men dens Mulighed, „Kraften“, der bestaar, og medens der
ikke bestandig er samme Sum af Bevægelse i Universet, er der samme
Mængde af „Kraft“ (Energi, efter det moderne Udtryk). Leibniz
lagde saa stor Vægt paa denne Sætning, at han endog kaldte den
Grundlaget for alle Naturens Love (\textit{fundamentum naturæ legum}). I
den følgende Tids Naturvidenskab traadte af forskellige Grunde
denne Sætning tilbage, skønt den atter og atter har virket som en
stiltiende Forudsætning. Da den nu kom frem igen, var det ikke som
en Sætning i den rene Mekanik, men som en Sætning, der hævdede
et bestemt Omsætningsforhold mellem de forskellige specielle Naturkræfter
(Bevægelse, Varme, Lys, Elektricitet o.\,s.\,v.), og den blev
nu gennemført ved experimentale Undersøgelser. Den staar nu som
en Lov for alle materielle Kraftytringer. --- Paa intet Sted viser det
sig saa tydelig som her, at det 19. Aarhundredes Videnskab i al sin
Storhed dog bestaar i Gennemførelse af Tanker, der stammer fra
den mærkelige Tid, da moderne Videnskab blev grundlagt.

Om Sætningen om Energiens Bestaaen er det blevet sagt, at den
kun er en Generalisation ud fra en begrænset Erfaring. Det er ogsaa
rigtigt, at en Erfaringsbekræftelse af en Sætning som denne i en vis
Forstand aldrig kan afsluttes. Man vil aldrig kunne isolere og afslutte
en Del af Verden saaledes, at enhver Tilførsel eller Udførsel
af Energi skulde være absolut udelukket. Og der staar stadig den
Mulighed tilbage, at Sætningen kun gælder for den Del eller de Dele
af Verden, man hidtil har kunnet undersøge. Og da „Verden“ aldrig
kan komme til at foreligge som et en Gang for alle afsluttet Hele,
% --- p. 29 ---
% running head: ROBERT MAYER OG SADI CARNOT (the page turns inside the word below)
er det en uendelig Opgave, der her er stillet Forskningen. Men Sæt\opage{29}ningen
om Energiens Bestaaen deler her kun Skæbne med alle de
Sætninger, Forskningen over den virkelige Natur kan føre til. Begrebet
„Verden“ i Betydning af absolut Totalitet har den videnskabelige
Betydning, at det viser os Muligheden af stedse nye Opgaver,
alt eftersom Horizonten udvides. Og det betegner et betydningsfuldt
Fremskridt, at den i det 17. Aarhundrede opstillede Sætning nu har
vist sig tilgængelig for fremskridende Erfaringsbekræftelse. Som enhver
stor Hypotese afgiver Sætningen en vigtig Grundsætning for
Forskningen, den nemlig, stedse at søge efter nye Energiformer som
Ækvivalenter for dem, der forsvandt, og at søge efter de forudgaaende
Energiformer, af hvilke de nye er opstaaede. Maxwell har særlig
fremhævet den erkendelsesteoretiske Betydning af Sætningen
eller Principet om Energiens Bestaaen, at vi føres til ved hvert nyt
Naturfænomen at stille bestemte Spørgsmaal. Det er jo overhovedet
alle levende Principers Betydning i Videnskaben, og denne Betydning
kan ikke overdrives.

Sætningen siger, at der kun finder Omsætning af Energi Sted ved
enhver materiel Forandring, uden at der gaar Energi tabt, men den
siger intet om, i hvilken Retning Omsætningerne finder Sted, og
navnlig om, hvorvidt en Omsætning i omvendt Retning altid ligesaa
let kan finde Sted. En fransk Forsker havde allerede 18 Aar forud
for Mayer’s Afhandling opstillet en Sætning, der belyser dette Forhold.
Sadi Carnot (Réflexions sur la puissance motrice du feu. 1824)
paaviste nemlig, at for at Varme skal være en bevægende Kraft, maa
der være Temperaturforskel tilstede. Bevægelse frembringes kun,
naar der indtræder en saadan Forskel mellem de Legemers Tilstand,
mellem hvilke Paavirkningen skal finde Sted. Men da Temperaturen
stadig har en Tendens til Ligevægt, idet Varmen gaar over fra varmere
til koldere Legemer og saaledes fordeles, falder derved indenfor
samme Kreds af Legemer, betragtet som et afsluttet Hele, Muligheden
af ny Omsætning af Varme til Bevægelse bort.

Ogsaa Carnot’s Sætning gælder kun for en begrænset Verden. Og
hvis man ved „Verden“ forstaar en begrænset Helhed, raader der i
den en Proces, som fører i Retning af en fuldstændig Ligevægt, en
Tilstand, i hvilken ingen ny Udløsning af Energi er mulig.

Carnot’s Sætning, der ligesom Mayer’s, maatte vente længe paa
fuld Anerkendelse, strider ikke mod Sætningen om Energiens Bestaaen.
Ti i den Ligevægtstilstand, der er stadig Tendens til at frembringe,
er ingen Energi gaaet tabt, kun er al Energi gaaet over fra
at være aktuel til at være potentiel.

% --- p. 30 ---
% running head: HORIZONTEN UDVIDES NEDAD OG UDAD
\opage{30}En optimistisk Naturfilosofi vil holde sig til Mayer’s Sætning, en
pessimistisk til Carnot’s. Men begge Sætninger gælder kun for en
begrænset Horizont, og da Naturvidenskaben stedse søger at udvide
sin Horizont, vil Diskussionen om, hvilket Synspunkt der er det afgørende,
stedse være at optage paany. ---

En endnu inderligere Sammenhæng mellem de forskellige Naturkræfter
end den, der følger af Sætningen om Energiens Bestaaen,
paavistes, da det blev godtgjort, at Naturkræfter, der hidtil stod som
forskelligartede, følger samme Love og derfor maa antages at være
af samme Natur. Saaledes paaviste Melloni og andre Forskere, at
Lys og Varme var Et. Allerede H. C. Ørsted anede en Sammenhæng
mellem Lys og Elektricitet, idet han opfattede Lynstraalen som
en Række elektriske Gnister. Og efterat Maxwell ad matematisk Vej
havde opstillet den elektromagnetiske Lysteori, ifølge hvilken Lys
kan føres tilbage til Elektricitet, idet Lovene for disse to Virksomheder
er de samme, paaviste Hertz experimentalt, at Lys og Elektricitet,
hvad Bølgebevægelse, Hastighed, Brydning o.\,s.\,v. angaar,
frembyder de samme Forhold. Optik kunde nu betragtes som et Tillæg
til Elektricitetslæren.

Overhovedet er det langt fra, at den Erfaringsbekræftelse, Sætningen
om Stoffets og Kraftens Bestaaen har vundet i det 19. Aarhundrede,
har bragt den naturvidenskabelige Verdensopfattelse til
Afslutning. Det var i sig selv et storartet Verdensbillede, der oprulledes.
Verden kom til at staa som en Helhed, indenfor hvilken stadige
Stof- og Kraftomsætninger fandt Sted, uden at der opstod eller forgik
Stof eller Kraft idethele. Som vi allerede har set, kan dette Verdensbillede
ikke afsluttes, da Horizontens Udvidelse stedse kan
bringe nye Muligheder. Men ogsaa paa en anden Kant, i selve Grundlaget,
gaar Forskningen ud over det antydede Verdensbillede. Elektricitetslærens
Udvikling har ført til, at der, langt nede under tidligere
Videnskabs Stoffer og Kræfter, dukker en ny og rig Verden
frem, i hvilken maaske det egentlige Grundlag for den Stoffets og
Kraftens Verden, i hvilken vi færdes, er at søge. Elektricitetslæren
har ført til Opdagelse af Smaadele, saakaldte Elektroner, der er
mindre end Dalton’s Atomer, og af hvilke disse maa tænkes sammensatte.
Den Forandring, der derved ved Aarhundredets Slutning
indvarsledes i den naturvidenskabelige Opfattelse, beskrives paa følgende
Maade af en dansk Fysiker. „I den sidste Halvdel af Aarhundredet
har Elektricitetslæren trængt sig saa stærkt frem paa de andre
Naturkræfters Bekostning, at den synes at ville opsluge dem
% --- p. 31 ---
% running head: EN NY FYSIK
\opage{31}alle. Man er nu paa Forhaand tilbøjelig til at henføre ikke alene
Lyslæren, men ogsaa de andre Naturkræfter til den; Elektriciteten
alene er saa mangfoldig i sine Ytringsformer, at den kan tænkes at
kunne omfatte dem alle. Endogsaa selve Begrebet Masse kan undværes,
naar Atomerne tænkes at være ladede med Elektricitet, idet
enhver Acceleration i saa Fald maa fremkalde Induktion i Omgivelserne
og derved igen en Modstand mod Bevægelsen. Elektricitetslærens
opløsende Kraft har ogsaa udstrakt sig til Kemien, hvorom
den fysiske Kemi bærer Vidnesbyrd. Selv det gamle Atombegreb
staar i Fare for at styrte sammen. Saaledes har J. J. Thomson og
Zeemann ment at kunne paavise Masser, der er tusinde Gange lettere
% printed thus: »Zeemann« for »Zeeman«; transcribed as printed
end Kemikernes Atomer“ (C. Christiansen: „Før og nu“, første
Bind af Tidsskrift for Fysik. 1908).

Der aabner sig saaledes ved Aarhundredets Slutning Udsigt til en
hel ny Fysik, som gaar længere tilbage end den af Galilei og Newton
grundede og vil aflede den gamle Fysiks Grundsætninger af endnu
dybere liggende Principer. Fra de store Massers Verden er man begyndt
at stige ned til de smaa Elementers Verden. „Den Tid,“ siger
en amerikansk Fysiker, „synes ikke at være fjern, da der vil blive
Plads for nye „Principia“ og for et Værk, der vil fuldføre for de
molekulære Systemer, hvad „Mécanique Céliste“ fuldførte for Solsystemet“.
% printed thus: »Céliste« for »Céleste« (Laplace, Mécanique céleste); transcribed as printed
(Woodward: The Unity of Physical Science. Congress of
Arts and Sciences St. Louis 1904).

Spørges der, hvad det da er, der bevæger sig efter Elektricitetens
Love og under specielle Betingelser fremtræder som Lys, Varme
o.\,s.\,v., svarer Naturforskeren med det smukke Navn Æter. Men
Æteren kendes kun gennem sine specielle Ytringsformer, blandt
hvilke Elektriciteten betragtes som den fundamentale. Den berømte
franske Fysiker Langevin erklærer, at det er umuligt at give et Billede,
en Forestilling om Æteren, da dens Egenskaber er helt forskellige
fra Materiens og sandsynligvis langt simplere end dennes.
Og han mener, at den Vanskelighed, vi her staar overfor, maaske
hænger sammen med den menneskelige Bevidstheds nuværende Beskaffenhed:
vi har i Aarhundreder været vante at tænke ved Hjælp
af Materiebegrebet og er derfor ude af Stand til at danne en
konkret Forestilling, som ikke har et materielt Indhold. Æteren
kan kun tænkes ved Hjælp af Begreberne om de elektromagnetiske
Egenskaber, vi tillægger den (Langevin: Congress of Arts and
Sciences. 1904).

% --- p. 32 ---
% running head: MEKANISME OG VITALISME (the page turns inside the word below)
Allerede Platon saa, at en gennemført Atomlære ikke kan opbyg\opage{32}ges
ved Hjælp af det sædvanlige Materiebegreb, ikke grundes paa
sanselige Forestillinger, men maa benytte sig af Analogi. Det Resultat,
han naaede til ved kritisk Gennemtænkning af Demokrits Atomlære,
er den moderne Fysik nu ad sin egen Vej naaet til. ---

For ingen Videnskab har der ved det tyvende Aarhundredes Begyndelse
vist sig større Muligheder og Opgaver end for Fysiken. Men
allerede, hvad der er vundet, --- fremskridende Gennemførelse og
Bekræftelse af de store Tanker, der for tre Hundrede Aar siden førte
til en „ny Videnskabs“ Opstaaen, er en stor Bedrift, der vil blive af
uoverskuelig Betydning for den menneskelige Opfattelse af Tilværelsen,
naar dens Rækkevidde engang ret gaar op for Bevidstheden.
Baade i det Store og i det Smaa, i Himmellegemernes og i de uanskuelige
Elementers Verden har der vist sig en større Sammenhæng
end man kunde ane.

\sectionhead{b.}{Udviklingslæren.}

Videnskaben om det organiske Liv kunde ikke forblive uberørt
af den store Bevægelse, der rørte sig paa det fysiske og kemiske
Omraade. Lige fra den nyere Naturvidenskabs første Opstaaen havde
man søgt at indordne Læren om Livet under de almindelige Naturlove,
saaledes som Fysiken paaviser dem. Den Sætning, at al Forandring
er Bevægelse, søgte man at hævde ikke blot paa den uorganiske,
men ogsaa paa den organiske Naturs Omraade. Descartes
gjorde med Eftertryk gældende, at Organismerne kun var Mekanismer,
der sattes i Gang efter Bevægelsens almindelige Love, og han
mente, at Harvey’s Opdagelse af Blodets Kredsløb bekræftede denne
Opfattelse, der blev den herskende i det 17. Aarhundrede. Som Reaktion
mod denne mekaniske Opfattelse, optraadte den saakaldte
Vitalisme, der hævdede, at de organiske Fænomener kun kunde finde
deres Forklaring, naar man antog en særegen, af alle fysiske og kemiske
Kræfter uafhængig Livskraft. Sin mest udprægede Form
naaede Vitalismen i Lægeskolen i Montpellier i det 18. Aarhundrede.
„Livskraften“ leder de fysiske og kemiske Processer i Organismen i
bestemte Retninger, bevirker en Sammenhæng mellem alt, hvad der
forgaar i det levende Væsen, saa at dette kan bestaa som Helhed,
sørger for Dannelsen af de Stoffer, Livet behøver, og for tilstrækkelig
Styrke til at overvinde den Modstand, Livet møder udefra.

Det 19. Aarhundrede fremtræder her, ligesom paa andre Punkter,
som en Gennemførelse og Specialisering af det 17. Aarhundredes
Grundtanker. Mekanisme og Vitalisme betegner i Grunden kun to
[\dots]

\end{document}
